% GENERATED FILE. Edit canonical lessons, then run npm run generate:books.
% canonical-source-hash: fnv1a64:35da96e609471416
% canonical-lessons: IT-C353-reazione, IT-C353-reagire, IT-C353-pretesto, IT-C353-scusa, IT-C353-giustificare

\chapter{Come reagire}
\label{ch:it-come-reagire}
\hlchaptermodality{353}

\begin{chapteropening}
\textbf{By the end of this chapter:} \emph{I can describe how someone reacted, and tell an honest reason from a pretext or an excuse and justify myself.}

Complete the last lesson of chapter 353: I can describe how someone reacted, and tell an honest reason from a pretext or an excuse and justify myself.
\end{chapteropening}

\section[la reazione]{la reazione — the reaction}

\label{lesson:IT-C353-reazione}



\begin{quote}
Before the new one: say the Italian for to unleash, then the Italian for to trigger.
\end{quote}

\subsection*{You'll want to know: la reazione}
\textbf{la reazione} — \textquotedblleft{}the reaction\textquotedblright{}.

\begin{quote}
La sua reazione mi ha sorpreso.
\end{quote}

\begin{quote}
\emph{Her reaction surprised me.}
\end{quote}

\emph{Una reazione forte} — a strong reaction.

\begin{cousinweb}
\emph{reazione} is built from \emph{re-} \textquotedblleft{}back\textquotedblright{} and \emph{azione} \textquotedblleft{}action\textquotedblright{}, a learned borrowing from Latin \emph{āctiō} \textquotedblleft{}a doing\textquotedblright{}, from \emph{agere} \textquotedblleft{}to do, to drive\textquotedblright{} (Lewis \& Short).
\end{cousinweb}

\begin{grammarlens}[title={What you've built}]
One more word for telling, explaining and describing at B1.
\end{grammarlens}

\subsection*{Your turn}
\begin{itemize}
\raggedright
  \item \emph{Say it:} \emph{la reazione}
  \item \emph{Say it:} \emph{la reazione}, once more
  \item \emph{Say it:} a sentence of your own with \emph{la reazione}
  \item \emph{Recall:} say the Italian for to unleash, then the Italian for to trigger, then say \emph{la reazione} again
\end{itemize}

\begin{tcolorbox}[breakable,colback=teal!4,colframe=teal!35!black,title={Before you move on}]
What does la reazione mean? (\textquotedblleft{}The reaction\textquotedblright{}.) Say it once more.
\end{tcolorbox}
\section[reagire]{reagire — to react}

\label{lesson:IT-C353-reagire}



\begin{quote}
Before the new one: say the Italian for to trigger, then the Italian for the reaction.
\end{quote}

\subsection*{You'll want to know: reagire}
\textbf{reagire} — \textquotedblleft{}to react\textquotedblright{}.

\begin{quote}
Come hai reagito quando hai sentito la notizia?
\end{quote}

\begin{quote}
\emph{How did you react when you heard the news?}
\end{quote}

Conjugated with \emph{-isc-}: \emph{reagisco, reagisci}.

\begin{cousinweb}
\emph{reagire} is built from \emph{re-} \textquotedblleft{}back\textquotedblright{} and \emph{agire} \textquotedblleft{}to act\textquotedblright{}, a learned borrowing from Latin \emph{agere} \textquotedblleft{}to do, to drive\textquotedblright{} (Lewis \& Short).
\end{cousinweb}

\begin{grammarlens}[title={What you've built}]
One more word for telling, explaining and describing at B1.
\end{grammarlens}

\subsection*{Your turn}
\begin{itemize}
\raggedright
  \item \emph{Say it:} \emph{reagire}
  \item \emph{Say it:} \emph{reagire}, once more
  \item \emph{Say it:} a sentence of your own with \emph{reagire}
  \item \emph{Recall:} say the Italian for to trigger, then the Italian for the reaction, then say \emph{reagire} again
\end{itemize}

\begin{tcolorbox}[breakable,colback=teal!4,colframe=teal!35!black,title={Before you move on}]
What does reagire mean? (\textquotedblleft{}To react\textquotedblright{}.) Say it once more.
\end{tcolorbox}
\section[il pretesto]{il pretesto — the pretext}

\label{lesson:IT-C353-pretesto}



\begin{quote}
Before the new one: say the Italian for the reaction, then the Italian for to react.
\end{quote}

\subsection*{You'll want to know: il pretesto}
\textbf{il pretesto} — \textquotedblleft{}the pretext\textquotedblright{}.

\begin{quote}
Ha usato il mal di testa come pretesto per non uscire.
\end{quote}

\begin{quote}
\emph{He used the headache as a pretext for not going out.}
\end{quote}

\emph{Con il pretesto di} — on the pretext of.

\begin{cousinweb}
\emph{pretesto} is a learned borrowing from Latin \emph{praetextus} \textquotedblleft{}an outward show, a pretext\textquotedblright{} (Lewis \& Short), from \emph{praetexere} \textquotedblleft{}to weave in front, to cover\textquotedblright{}, built on \emph{texere} \textquotedblleft{}to weave\textquotedblright{}.
\end{cousinweb}

\begin{grammarlens}[title={What you've built}]
One more word for telling, explaining and describing at B1.
\end{grammarlens}

\subsection*{Your turn}
\begin{itemize}
\raggedright
  \item \emph{Say it:} \emph{il pretesto}
  \item \emph{Say it:} \emph{il pretesto}, once more
  \item \emph{Say it:} a sentence of your own with \emph{il pretesto}
  \item \emph{Recall:} say the Italian for the reaction, then the Italian for to react, then say \emph{il pretesto} again
\end{itemize}

\begin{tcolorbox}[breakable,colback=teal!4,colframe=teal!35!black,title={Before you move on}]
What does il pretesto mean? (\textquotedblleft{}The pretext\textquotedblright{}.) Say it once more.
\end{tcolorbox}
\section[la scusa]{la scusa — the excuse; the apology}

\label{lesson:IT-C353-scusa}



\begin{quote}
Before the new one: say the Italian for to react, then the Italian for the pretext.
\end{quote}

\subsection*{You'll want to know: la scusa}
\textbf{la scusa} — \textquotedblleft{}the excuse; the apology\textquotedblright{}.

\begin{quote}
Non cercare scuse: hai sbagliato tu.
\end{quote}

\begin{quote}
\emph{Do not look for excuses: you got it wrong.}
\end{quote}

\emph{Chiedere scusa} — to apologise. \emph{Scusa!} — sorry! (to a friend); \emph{scusi!} to a stranger.

\begin{cousinweb}
\emph{scusa} is built on \emph{scusare} \textquotedblleft{}to excuse\textquotedblright{}, inherited from Latin \emph{excūsāre} \textquotedblleft{}to free from a charge\textquotedblright{} (Lewis \& Short), built from \emph{ex} \textquotedblleft{}out of\textquotedblright{} and \emph{causa} \textquotedblleft{}a case, a charge\textquotedblright{}.
\end{cousinweb}

\begin{grammarlens}[title={What you've built}]
One more word for telling, explaining and describing at B1.
\end{grammarlens}

\subsection*{Your turn}
\begin{itemize}
\raggedright
  \item \emph{Say it:} \emph{la scusa}
  \item \emph{Say it:} \emph{la scusa}, once more
  \item \emph{Say it:} a sentence of your own with \emph{la scusa}
  \item \emph{Recall:} say the Italian for to react, then the Italian for the pretext, then say \emph{la scusa} again
\end{itemize}

\begin{tcolorbox}[breakable,colback=teal!4,colframe=teal!35!black,title={Before you move on}]
What does la scusa mean? (\textquotedblleft{}The excuse; the apology\textquotedblright{}.) Say it once more.
\end{tcolorbox}
\section[giustificare]{giustificare — to justify}

\label{lesson:IT-C353-giustificare}



\begin{quote}
Before the new one: say the Italian for the pretext, then the Italian for the excuse.
\end{quote}

\subsection*{You'll want to know: giustificare}
\textbf{giustificare} — \textquotedblleft{}to justify\textquotedblright{}.

\begin{quote}
Niente può giustificare un comportamento così.
\end{quote}

\begin{quote}
\emph{Nothing can justify behaviour like that.}
\end{quote}

\emph{Giustificarsi} — to justify yourself. At school, \emph{giustificare un ritardo} — to give a reason for being late.

\begin{cousinweb}
\emph{giustificare} is a learned borrowing from Late Latin \emph{iūstificāre} \textquotedblleft{}to deal justly with, to justify\textquotedblright{} (Lewis \& Short), built from \emph{iūstus} \textquotedblleft{}right, fair\textquotedblright{} and \emph{facere} \textquotedblleft{}to make\textquotedblright{}.
\end{cousinweb}

\begin{grammarlens}[title={What you've built}]
One more word for telling, explaining and describing at B1.
\end{grammarlens}

\subsection*{Your turn}
\begin{itemize}
\raggedright
  \item \emph{Say it:} \emph{giustificare}
  \item \emph{Say it:} \emph{giustificare}, once more
  \item \emph{Say it:} a sentence of your own with \emph{giustificare}
  \item \emph{Recall:} say the Italian for the pretext, then the Italian for the excuse, then say \emph{giustificare} again
\end{itemize}

\begin{tcolorbox}[breakable,colback=teal!4,colframe=teal!35!black,title={Before you move on}]
What does giustificare mean? (\textquotedblleft{}To justify\textquotedblright{}.) Say it once more.
\end{tcolorbox}
